% Jesús del 1 - 9

% Hacer reserva de actividad
% Alta cliente
% Listar reservas
% Pagar reserva
% Pago con PayPal
% Pago con Bizum
% Pago con tarjeta
% Cancelar reserva
% Ver mis reservas


\chapter{Gestión de Reservas}

% ======Tabla Caso de uso Número: 1 =========================
% Please add the following required packages to your document preamble:
% \usepackage[table,xcdraw]{xcolor}
% Beamer presentation requires \usepackage{colortbl} instead of \usepackage[table,xcdraw]{xcolor}
\begin{table}[H]
    \begin{tabular}{llllllll}
    \hline
    \multicolumn{1}{|l|}{Caso de Uso}                            & \multicolumn{5}{l|}{{\color[HTML]{808080} \textit{Hacer reserva de actividad}}}                                                                                                                                                                & \multicolumn{2}{l|}{\cellcolor[HTML]{E6E6E6}{\color[HTML]{808080} \textit{CU-1}}}                                                   \\ \hline
    \multicolumn{1}{|l|}{Actores}                                & \multicolumn{7}{l|}{{\color[HTML]{808080} \textit{Cliente(I), Cliente Registrado}}}                                                                                                                                                                                                                                                                                               \\ \hline
    \multicolumn{1}{|l|}{Tipo}                                   & \multicolumn{7}{l|}{{\color[HTML]{808080} \textit{Primario Esencial}}}                                                                                                                                                                                                                                                                                                             \\ \hline
    \multicolumn{1}{|l|}{Referencias}                            & \multicolumn{2}{l|}{}                                                                                                   & \multicolumn{5}{l|}{\color[HTML]{808080}\textit{CU-4 Pagar reserva (incluye)}}                                                                                                                                                                                                                                    \\ \hline
    \multicolumn{1}{|l|}{Precondición}                           & \multicolumn{7}{p{10cm}|}{{\color[HTML]{808080} \textit{}}}                                                                                                                                                                                                                                   \\ \hline
    \multicolumn{1}{|l|}{Poscondición}                           & \multicolumn{7}{l|}{{\color[HTML]{808080} \textit{}}}                                                                                                                                                                                                                                                       \\ \hline
    \multicolumn{1}{|l|}{Autor}                                  & \multicolumn{1}{l|}{{\color[HTML]{808080} \textit{Jesús Rodríguez González}}}                           & \multicolumn{2}{l|}{Fecha}                                                   & \multicolumn{1}{l|}{}                          & \multicolumn{2}{l|}{Versión}                                                      & \multicolumn{1}{l|}{1.0}                                                     \\ \hline
    \end{tabular}
    \end{table}

    \begin{table}[H]
        \begin{tabular}{llllllll}
    \hline
    \multicolumn{8}{|l|}{\cellcolor[HTML]{F2F2F2}Propósito}                                                                                                                                                                                                                                                                                                                                                                                           \\ \hline
    \multicolumn{8}{|p{15cm}|}{{\color[HTML]{808080} \textit{Permitir a un Cliente (registrado o no) realizar la reserva de una actividad.}}}                                                                                                                                                                                                             \\ \hline
    \end{tabular}
    \end{table}

    \begin{table}[H]
        \begin{tabular}{llllllll}
    \hline
    \multicolumn{8}{|l|}{\cellcolor[HTML]{F2F2F2}Resumen}                                                                                                                                                                                                                                                                                                                                                                                             \\ \hline
    \multicolumn{8}{|p{15cm}|}{{\color[HTML]{808080} \textit{Un Cliente puede decidir registrarse o no en la aplicación para realizar la reserva de una aplicación. Decida o no darse de alta tras comenzar a hacer una reserva, es obligatorio que se realice el pago en alguna de sus modalidades disponibles.}}} \\ \hline
    \end{tabular}
    \end{table}


% Please add the following required packages to your document preamble:
% \usepackage[table,xcdraw]{xcolor}
% Beamer presentation requires \usepackage{colortbl} instead of \usepackage[table,xcdraw]{xcolor}

\definecolor{lightgray}{gray}{0.9}

\begin{table}[h!]
    \centering
    \begin{tabular}{|p{0.6cm}|p{6.5cm}|p{0.6cm}|p{6.5cm}|}
    \rowcolor{lightgray}
    \hline
    \multicolumn{4}{|l|}{Curso Normal (Básico)} \\
    \hline
    1 & \color[HTML]{808080}\textit{Un cliente solicita reservar una actividad un día a una hora concretos.} & 2 & \color[HTML]{808080}\textit{El sistema comprueba la disponibilidad de la actividad en la fecha solicitada.} \\
    \hline
    & & 3 & \color[HTML]{808080}\textit{El sistema pide al usuario sus datos.} \\
    \hline
    & \color[HTML]{808080}\textit{Punto de extensión: Alta cliente} & & \\
    \hline
    4 & \color[HTML]{808080}\textit{El cliente introduce sus datos para la reserva.} & 5 & \color[HTML]{808080}\textit{El sistema comprueba si los datos son válidos.} \\
    \hline
    6 & \color[HTML]{808080}\textit{El cliente procede al pago de la reserva} & 7 & \color[HTML]{808080}\textit{Incluir CU-4 (Pagar reserva)} \\
    \hline
    & & 8 & \color[HTML]{808080}\textit{Si el pago es correcto, el sistema almacena la reserva y envía un justificante al cliente.} \\
    \hline
    \end{tabular}
    \end{table}

    \begin{table}[h!]
        \centering
        \begin{tabular}{|p{0.6cm}|p{13.4cm}|}
        \rowcolor{lightgray}
        \hline
        \multicolumn{2}{|l|}{Cursos Alternos} \\
        \hline
        2a & \color[HTML]{808080}\textit{Si la actividad no está disponible, se notifica al cliente y termina el caso de uso.} \\
        \hline
        5a & \color[HTML]{808080}\textit{Si los datos no son válidos, se notifica al cliente y termina el caso de uso.} \\
        \hline
        8a & \color[HTML]{808080}\textit{Si no se ha podido realizar el pago correctamente, se notifica al cliente.} \\
        \hline
        \end{tabular}
        \end{table}
        
        
        
    % Please add the following required packages to your document preamble:
    % \usepackage[table,xcdraw]{xcolor}
    % Beamer presentation requires \usepackage{colortbl} instead of \usepackage[table,xcdraw]{xcolor}
    \begin{table}[H]
    \begin{tabular}{|l
    >{\columncolor[HTML]{FFFFFF}}l l
    >{\columncolor[HTML]{FFFFFF}}l 
    >{\columncolor[HTML]{FFFFFF}}l 
    >{\columncolor[HTML]{FFFFFF}}l 
    >{\columncolor[HTML]{FFFFFF}}l 
    >{\columncolor[HTML]{FFFFFF}}l |}
    \hline
    \multicolumn{8}{|l|}{\cellcolor[HTML]{F2F2F2}Otros datos}                                                                                                                                                                                                                                                                                                        \\ \hline
    \multicolumn{1}{|l|}{Frecuencia esperada} & \multicolumn{1}{p{5cm}|}{\cellcolor[HTML]{FFFFFF}{\color[HTML]{808080} }} & \multicolumn{1}{l|}{Rendimiento} & \multicolumn{5}{p{5cm}|}{\cellcolor[HTML]{FFFFFF}{\color[HTML]{808080} }} \\ \hline
    \multicolumn{1}{|l|}{Importancia}         & \multicolumn{1}{p{5cm}|}{\cellcolor[HTML]{FFFFFF}{\color[HTML]{808080} }}   & \multicolumn{1}{l|}{Urgencia}    & \multicolumn{5}{p{5cm}|}{\cellcolor[HTML]{FFFFFF}{\color[HTML]{808080} }}                          \\ \hline
    \multicolumn{1}{|l|}{Estado}              & \multicolumn{1}{p{5cm}|}{\cellcolor[HTML]{FFFFFF}{\color[HTML]{808080}}}                   & \multicolumn{1}{l|}{Estabilidad} & \multicolumn{5}{p{5cm}|}{\cellcolor[HTML]{FFFFFF}{\color[HTML]{808080} }}                                            \\ \hline
    \end{tabular}
    \end{table}

    \begin{table}[h!]
        \centering
        \begin{tabular}{|p{14.2cm}|}
        \rowcolor{lightgray}
        \hline
        Comentarios \\
        \hline
         \\
        \hline
        \end{tabular}
        \end{table}

% ======Tabla Caso de uso Número: 2 =========================
\begin{table}[H]
    \begin{tabular}{llllllll}
    \hline
    \multicolumn{1}{|l|}{Caso de Uso}                            & \multicolumn{5}{l|}{{\color[HTML]{808080} \textit{Alta cliente}}}                                                                                                                                                                & \multicolumn{2}{l|}{\cellcolor[HTML]{E6E6E6}{\color[HTML]{808080} \textit{CU-2}}}                                                   \\ \hline
    \multicolumn{1}{|l|}{Actores}                                & \multicolumn{7}{l|}{{\color[HTML]{808080} \textit{Cliente (I)}}}                                                                                                                                                                                                                                                                                               \\ \hline
    \multicolumn{1}{|l|}{Tipo}                                   & \multicolumn{7}{l|}{{\color[HTML]{808080} \textit{Primario Esencial}}}                                                                                                                                                                                                                                                                                                             \\ \hline
    \multicolumn{1}{|l|}{Referencias}                            & \multicolumn{2}{l|}{}                                                                                                   & \multicolumn{5}{l|}{\color[HTML]{808080} \textit{CU-1 Hacer reserva de actividad (extiende)}}                                                                                                                                                                                                                                    \\ \hline
    \multicolumn{1}{|l|}{Precondición}                           & \multicolumn{7}{p{10cm}|}{{\color[HTML]{808080} \textit{El cliente no debe estar registrado en el sistema.}}}                                                                                                                                                                                                                                   \\ \hline
    \multicolumn{1}{|l|}{Poscondición}                           & \multicolumn{7}{l|}{{\color[HTML]{808080} \textit{El cliente queda registrado en el sistema.}}}                                                                                                                                                                                                                                                       \\ \hline
    \multicolumn{1}{|l|}{Autor}                                  & \multicolumn{1}{l|}{{\color[HTML]{808080} \textit{Jesús Rodríguez González}}}                           & \multicolumn{2}{l|}{Fecha}                                                   & \multicolumn{1}{l|}{}                          & \multicolumn{2}{l|}{Versión}                                                      & \multicolumn{1}{l|}{1.0}                                                     \\ \hline
    \end{tabular}
    \end{table}

    \begin{table}[H]
        \begin{tabular}{llllllll}
    \hline
    \multicolumn{8}{|l|}{\cellcolor[HTML]{F2F2F2}Propósito}                                                                                                                                                                                                                                                                                                                                                                                           \\ \hline
    \multicolumn{8}{|p{15cm}|}{{\color[HTML]{808080} \textit{Permitir a un cliente no registrado darse de alta en la aplicación.}}}                                                                                                                                                                                                             \\ \hline
    \end{tabular}
    \end{table}

    \begin{table}[H]
        \begin{tabular}{llllllll}
    \hline
    \multicolumn{8}{|l|}{\cellcolor[HTML]{F2F2F2}Resumen}                                                                                                                                                                                                                                                                                                                                                                                             \\ \hline
    \multicolumn{8}{|p{15cm}|}{{\color[HTML]{808080} \textit{Este caso de uso permite que un cliente se registre en la aplicación, lo
    cual le autoriza a realizar una serie de funcionalidades sobre sus reservas. Para ello deberá introducir sus datos personales y que sean aceptados por el sistema.}}} \\ \hline
    \end{tabular}
    \end{table}

\begin{table}[h!]
    \centering
    \begin{tabular}{|p{0.6cm}|p{6.5cm}|p{0.6cm}|p{6.5cm}|}
    \rowcolor{lightgray}
    \hline
    \multicolumn{4}{|l|}{Curso Normal (Básico)} \\
    \hline
    1 & \color[HTML]{808080}\textit{Un cliente se quiere dar de alta en la aplicación.} & 2 & \color[HTML]{808080}\textit{El sistema pide al usuario que introduzca sus datos.} \\
    \hline
    3 & \color[HTML]{808080}\textit{El cliente introduce sus datos personales.} & 4 & \color[HTML]{808080}\textit{El sistema comprueba que los datos introducidos sean válidos.} \\
    \hline
    & & 5 & \color[HTML]{808080}\textit{El sistema comprueba que no exista un usuario registrado con esos datos.} \\
    \hline
    & & 6 & \color[HTML]{808080}\textit{El sistema notifica al usuario que se ha registrado con éxito.} \\
    \hline
    \end{tabular}
    \end{table}

    \begin{table}[h!]
        \centering
        \begin{tabular}{|p{0.6cm}|p{13.4cm}|}
        \rowcolor{lightgray}
        \hline
        \multicolumn{2}{|l|}{Cursos Alternos} \\
        \hline
        4a & \color[HTML]{808080}\textit{Si los datos introducidos no son válidos, se notifica al cliente y termina el caso de uso.} \\
        \hline
        5a & \color[HTML]{808080}\textit{Si el usuario ya existía, se notifica al cliente y termina el caso de uso.} \\
        \hline

        \end{tabular}
        \end{table}
        


    % Please add the following required packages to your document preamble:
    % \usepackage[table,xcdraw]{xcolor}
    % Beamer presentation requires \usepackage{colortbl} instead of \usepackage[table,xcdraw]{xcolor}
    \begin{table}[H]
        \begin{tabular}{|l
        >{\columncolor[HTML]{FFFFFF}}l l
        >{\columncolor[HTML]{FFFFFF}}l 
        >{\columncolor[HTML]{FFFFFF}}l 
        >{\columncolor[HTML]{FFFFFF}}l 
        >{\columncolor[HTML]{FFFFFF}}l 
        >{\columncolor[HTML]{FFFFFF}}l |}
        \hline
        \multicolumn{8}{|l|}{\cellcolor[HTML]{F2F2F2}Otros datos}                                                                                                                                                                                                                                                                                                        \\ \hline
        \multicolumn{1}{|l|}{Frecuencia esperada} & \multicolumn{1}{p{5cm}|}{\cellcolor[HTML]{FFFFFF}{\color[HTML]{808080} }} & \multicolumn{1}{l|}{Rendimiento} & \multicolumn{5}{p{5cm}|}{\cellcolor[HTML]{FFFFFF}{\color[HTML]{808080} }} \\ \hline
        \multicolumn{1}{|l|}{Importancia}         & \multicolumn{1}{p{5cm}|}{\cellcolor[HTML]{FFFFFF}{\color[HTML]{808080} }}   & \multicolumn{1}{l|}{Urgencia}    & \multicolumn{5}{p{5cm}|}{\cellcolor[HTML]{FFFFFF}{\color[HTML]{808080} }}                          \\ \hline
        \multicolumn{1}{|l|}{Estado}              & \multicolumn{1}{p{5cm}|}{\cellcolor[HTML]{FFFFFF}{\color[HTML]{808080}}}                   & \multicolumn{1}{l|}{Estabilidad} & \multicolumn{5}{p{5cm}|}{\cellcolor[HTML]{FFFFFF}{\color[HTML]{808080} }}                                            \\ \hline
        \end{tabular}
        \end{table}

        \begin{table}[h!]
            \centering
            \begin{tabular}{|p{14.2cm}|}
            \rowcolor{lightgray}
            \hline
            Comentarios \\
            \hline
             \\
            \hline
            \end{tabular}
            \end{table}

    % ======Tabla Caso de uso Número: 3 =========================
    \begin{table}[H]
        \begin{tabular}{llllllll}
        \hline
        \multicolumn{1}{|l|}{Caso de Uso}                            & \multicolumn{5}{l|}{{\color[HTML]{808080} \textit{Listar reservas}}}                                                                                                                                                                & \multicolumn{2}{l|}{\cellcolor[HTML]{E6E6E6}{\color[HTML]{808080} \textit{CU-3}}}                                                   \\ \hline
        \multicolumn{1}{|l|}{Actores}                                & \multicolumn{7}{l|}{{\color[HTML]{808080} \textit{Empleado(I)}}}                                                                                                                                                                                                                                                                                               \\ \hline
        \multicolumn{1}{|l|}{Tipo}                                   & \multicolumn{7}{l|}{{\color[HTML]{808080} \textit{Primario Esencial}}}                                                                                                                                                                                                                                                                                                             \\ \hline
        \multicolumn{1}{|l|}{Referencias}                            & \multicolumn{2}{l|}{}                                                                                                   & \multicolumn{5}{l|}{}                                                                                                                                                                                                                                    \\ \hline
        \multicolumn{1}{|l|}{Precondición}                           & \multicolumn{7}{p{10cm}|}{{\color[HTML]{808080} \textit{Un cliente debe haber comenzado el proceso de reserva de una actividad.}}}                                                                                                                                                                                                                                   \\ \hline
        \multicolumn{1}{|l|}{Poscondición}                           & \multicolumn{7}{l|}{{\color[HTML]{808080} \textit{}}}                                                                                                                                                                                                                                                       \\ \hline
        \multicolumn{1}{|l|}{Autor}                                  & \multicolumn{1}{l|}{{\color[HTML]{808080} \textit{Jesús Rodríguez González}}}                           & \multicolumn{2}{l|}{Fecha}                                                   & \multicolumn{1}{l|}{}                          & \multicolumn{2}{l|}{Versión}                                                      & \multicolumn{1}{l|}{1.0}                                                     \\ \hline
        \end{tabular}
        \end{table}

        \begin{table}[H]
            \begin{tabular}{llllllll}
        \hline
        \multicolumn{8}{|l|}{\cellcolor[HTML]{F2F2F2}Propósito}                                                                                                                                                                                                                                                                                                                                                                                           \\ \hline
        \multicolumn{8}{|p{15cm}|}{{\color[HTML]{808080} \textit{Este caso de uso permite que un actor Empleado obtenga
        información sobre todas las reservas que tiene la empresa.}}}                                                                                                                                                                                                             \\ \hline
        \end{tabular}
        \end{table}

        \begin{table}[H]
            \begin{tabular}{llllllll}
        \hline
        \multicolumn{8}{|l|}{\cellcolor[HTML]{F2F2F2}Resumen}                                                                                                                                                                                                                                                                                                                                                                                             \\ \hline
        \multicolumn{8}{|p{15cm}|}{{\color[HTML]{808080} \textit{Un Empleado podrá acceder a todas las reservas que hay almaceanadas en el sistema. En caso de que haya se mostrarán al Empleado, y en caso de que no haya aparecerá la lista vacía.}}} \\ \hline
        \end{tabular}
        \end{table}

    
\begin{table}[h!]
    \centering
    \begin{tabular}{|p{0.6cm}|p{6.5cm}|p{0.6cm}|p{6.5cm}|}
    \rowcolor{lightgray}
    \hline
    \multicolumn{4}{|l|}{Curso Normal (Básico)} \\
    \hline
    1 & \color[HTML]{808080}\textit{Un empleado solicita listar las reservas con su número de empleado.} & 2 & \color[HTML]{808080}\textit{El sistema comprueba que el número de empleado es válido.} \\
    \hline
    & & 3 & \color[HTML]{808080}\textit{El sistema busca las reservas almacenadas en su base de datos.} \\
    \hline
    & & 4 & \color[HTML]{808080}\textit{El sistema muestra al empleado la lista de reservas almacenada.} \\
    \hline
    \end{tabular}
    \end{table}

    \begin{table}[h!]
    \   \centering
        \begin{tabular}{|p{0.6cm}|p{13.4cm}|}
        \rowcolor{lightgray}
        \hline
        \multicolumn{2}{|l|}{Cursos Alternos} \\
        \hline
        2b & \color[HTML]{808080}\textit{Si el sistema comprueba que el número de empleado no es válido, se notifica al empleado y termina el caso de uso.} \\
        \hline
        \end{tabular}
        \end{table}
  

        \begin{table}[H]
            \begin{tabular}{|l
            >{\columncolor[HTML]{FFFFFF}}l l
            >{\columncolor[HTML]{FFFFFF}}l 
            >{\columncolor[HTML]{FFFFFF}}l 
            >{\columncolor[HTML]{FFFFFF}}l 
            >{\columncolor[HTML]{FFFFFF}}l 
            >{\columncolor[HTML]{FFFFFF}}l |}
            \hline
            \multicolumn{8}{|l|}{\cellcolor[HTML]{F2F2F2}Otros datos}                                                                                                                                                                                                                                                                                                        \\ \hline
            \multicolumn{1}{|l|}{Frecuencia esperada} & \multicolumn{1}{p{5cm}|}{\cellcolor[HTML]{FFFFFF}{\color[HTML]{808080} }} & \multicolumn{1}{l|}{Rendimiento} & \multicolumn{5}{p{5cm}|}{\cellcolor[HTML]{FFFFFF}{\color[HTML]{808080} }} \\ \hline
            \multicolumn{1}{|l|}{Importancia}         & \multicolumn{1}{p{5cm}|}{\cellcolor[HTML]{FFFFFF}{\color[HTML]{808080} }}   & \multicolumn{1}{l|}{Urgencia}    & \multicolumn{5}{p{5cm}|}{\cellcolor[HTML]{FFFFFF}{\color[HTML]{808080} }}                          \\ \hline
            \multicolumn{1}{|l|}{Estado}              & \multicolumn{1}{p{5cm}|}{\cellcolor[HTML]{FFFFFF}{\color[HTML]{808080}}}                   & \multicolumn{1}{l|}{Estabilidad} & \multicolumn{5}{p{5cm}|}{\cellcolor[HTML]{FFFFFF}{\color[HTML]{808080} }}                                            \\ \hline
            \end{tabular}
            \end{table}
    
    
            \begin{table}[h!]
                \centering
                \begin{tabular}{|p{14.2cm}|}
                \rowcolor{lightgray}
                \hline
                Comentarios \\
                \hline
                 \\
                \hline
                \end{tabular}
                \end{table}
    
    % ======Tabla Caso de uso Número: 4 =========================
    \begin{table}[H]
        \begin{tabular}{llllllll}
        \hline
        \multicolumn{1}{|l|}{Caso de Uso}                            & \multicolumn{5}{l|}{{\color[HTML]{808080} \textit{Pagar reserva}}}                                                                                                                                                                & \multicolumn{2}{l|}{\cellcolor[HTML]{E6E6E6}{\color[HTML]{808080} \textit{CU-4}}}                                                   \\ \hline
        \multicolumn{1}{|l|}{Actores}                                & \multicolumn{7}{l|}{{\color[HTML]{808080} \textit{Cliente(I), 
        Cliente Registrado}}}                                                                                                                                                                                                                                                                                               \\ \hline
        \multicolumn{1}{|l|}{Tipo}                                   & \multicolumn{7}{l|}{{\color[HTML]{808080} \textit{Primario Esencial}}}                                                                                                                                                                                                                                                                                                             \\ \hline
        \multicolumn{1}{|l|}{Referencias}                            & \multicolumn{2}{l|}{} & \multicolumn{5}{l|}{\color[HTML]{808080}\textit{CU-5 Pago con PayPal (hijo), CU-6 Pago con Bizum (hijo),CU-7 Pago con tarjeta (hijo).}}                                                                                                                                                                                                                                    \\ \hline
        \multicolumn{1}{|l|}{Precondición}                           & \multicolumn{7}{p{10cm}|}{{\color[HTML]{808080} \textit{}}}                                                                                                                                                                                                                                   \\ \hline
        \multicolumn{1}{|l|}{Poscondición}                           & \multicolumn{7}{p{10cm}|}{{\color[HTML]{808080} \textit{}}}                                                                                                                                                                                                                                                       \\ \hline
        \multicolumn{1}{|l|}{Autor}                                  & \multicolumn{1}{l|}{{\color[HTML]{808080} \textit{Jesús Rodríguez González}}}                           & \multicolumn{2}{l|}{Fecha}                                                   & \multicolumn{1}{l|}{}                          & \multicolumn{2}{l|}{Versión}                                                      & \multicolumn{1}{l|}{1.0}                                                     \\ \hline
        \end{tabular}
        \end{table}

        \begin{table}[H]
            \begin{tabular}{llllllll}
        \hline
        \multicolumn{8}{|l|}{\cellcolor[HTML]{F2F2F2}Propósito}                                                                                                                                                                                                                                                                                                                                                                                           \\ \hline
        \multicolumn{8}{|p{15cm}|}{{\color[HTML]{808080} \textit{Permitir a un Cliente (registrado o no) pagar la reservaa de una actividad.}}}                                                                                                                                                                                                             \\ \hline
        \end{tabular}
        \end{table}

        \begin{table}[H]
            \begin{tabular}{llllllll}
        \hline
        \multicolumn{8}{|l|}{\cellcolor[HTML]{F2F2F2}Resumen}                                                                                                                                                                                                                                                                                                                                                                                             \\ \hline
        \multicolumn{8}{|p{15cm}|}{{\color[HTML]{808080} \textit{Este caso de uso permite pagar la reserva, el pago se puede realizar de distintas formas: con tarjeta (caso de uso Pago con Tarjeta), Bizum (caso de uso Pago con Bizum), PayPal (caso de uso Pago con PayPal). Una vez se realice una reserva de una actividad, será necesario que se realice el pago de la misma.}}} \\ \hline
        \end{tabular}
        \end{table}

        \begin{table}[h!]
        \centering
        \begin{tabular}{|p{0.6cm}|p{6.5cm}|p{0.6cm}|p{6.5cm}|}
        \rowcolor{lightgray}
        \hline
        \multicolumn{4}{|l|}{Curso Normal (Básico)} \\
        \hline
        1 & \color[HTML]{808080}\textit{Un cliente se dispone a pagar una reserva.} & 2 & \color[HTML]{808080}\textit{El sistema muestra al cliente los diferentes métodos de pago disponibles.} \\
        \hline
        3 & \color[HTML]{808080}\textit{El cliente selecciona el método de pago deseado.} & & \\
        \hline
        4 & \color[HTML]{808080}\textit{El cliente realiza el pago.} & 5 & \color[HTML]{808080}\textit{El sistema notifica si el pago ha sido realizado correctamente o no.} \\
        \hline
        \end{tabular}
        \end{table}

        \begin{table}[h!]
            \centering
            \begin{tabular}{|p{0.6cm}|p{13.4cm}|}
            \rowcolor{lightgray}
            \hline
            \multicolumn{2}{|l|}{Cursos Alternos} \\
            \hline
            4a & \color[HTML]{808080}\textit{Si el pago es con PayPal, incluir CU-5 (Pago con PayPal).} \\
            \hline
            4b & \color[HTML]{808080}\textit{Si el pago es con Bizum, incluir CU-6 (Pago con Bizum).} \\
            \hline
            4c & \color[HTML]{808080}\textit{Si el pago es con tarjeta, incluir CU-7 (Pago con tarjeta).} \\
            \hline
            5a & \color[HTML]{808080}\textit{Si el pago se realizó correctamente, se envía un justificante al cliente con la reserva.} \\
            \hline
            5b & \color[HTML]{808080}\textit{Si el pago falló, se cancela la reserva en proceso.} \\
            \hline
            \end{tabular}
            \end{table}
            

            % Please add the following required packages to your document preamble:
    % \usepackage[table,xcdraw]{xcolor}
    % Beamer presentation requires \usepackage{colortbl} instead of \usepackage[table,xcdraw]{xcolor}
    \begin{table}[H]
        \begin{tabular}{|l
        >{\columncolor[HTML]{FFFFFF}}l l
        >{\columncolor[HTML]{FFFFFF}}l 
        >{\columncolor[HTML]{FFFFFF}}l 
        >{\columncolor[HTML]{FFFFFF}}l 
        >{\columncolor[HTML]{FFFFFF}}l 
        >{\columncolor[HTML]{FFFFFF}}l |}
        \hline
        \multicolumn{8}{|l|}{\cellcolor[HTML]{F2F2F2}Otros datos}                                                                                                                                                                                                                                                                                                        \\ \hline
        \multicolumn{1}{|l|}{Frecuencia esperada} & \multicolumn{1}{p{5cm}|}{\cellcolor[HTML]{FFFFFF}{\color[HTML]{808080} }} & \multicolumn{1}{l|}{Rendimiento} & \multicolumn{5}{p{5cm}|}{\cellcolor[HTML]{FFFFFF}{\color[HTML]{808080} }} \\ \hline
        \multicolumn{1}{|l|}{Importancia}         & \multicolumn{1}{p{5cm}|}{\cellcolor[HTML]{FFFFFF}{\color[HTML]{808080} }}   & \multicolumn{1}{l|}{Urgencia}    & \multicolumn{5}{p{5cm}|}{\cellcolor[HTML]{FFFFFF}{\color[HTML]{808080} }}                          \\ \hline
        \multicolumn{1}{|l|}{Estado}              & \multicolumn{1}{p{5cm}|}{\cellcolor[HTML]{FFFFFF}{\color[HTML]{808080}}}                   & \multicolumn{1}{l|}{Estabilidad} & \multicolumn{5}{p{5cm}|}{\cellcolor[HTML]{FFFFFF}{\color[HTML]{808080} }}                                            \\ \hline
        \end{tabular}
        \end{table}

        \begin{table}[h!]
            \centering
            \begin{tabular}{|p{14.2cm}|}
            \rowcolor{lightgray}
            \hline
            Comentarios \\
            \hline
             \\
            \hline
            \end{tabular}
            \end{table}

    % ======Tabla Caso de uso Número: 5 =========================
    \begin{table}[H]
        \begin{tabular}{llllllll}
        \hline
        \multicolumn{1}{|l|}{Caso de Uso}                            & \multicolumn{5}{l|}{{\color[HTML]{808080} \textit{Pago con PayPal}}}                                                                                                                                                                & \multicolumn{2}{l|}{\cellcolor[HTML]{E6E6E6}{\color[HTML]{808080} \textit{CU-5}}}                                                   \\ \hline
        \multicolumn{1}{|l|}{Actores}                                & \multicolumn{7}{l|}{{\color[HTML]{808080} \textit{Cliente(I), Cliente Registrado, Sistema PayPal}}}                                                                                                                                                                                                                                                                                               \\ \hline
        \multicolumn{1}{|l|}{Tipo}                                   & \multicolumn{7}{l|}{{\color[HTML]{808080} \textit{Primario Esencial}}}                                                                                                                                                                                                                                                                                                             \\ \hline
        \multicolumn{1}{|l|}{Referencias}                            & \multicolumn{2}{l|}{}                                                                                                   & \multicolumn{5}{l|}{\color[HTML]{808080}\textit{CU-4 Pagar reserva (padre).}}                                                                                                                                                                                                                                   \\ \hline
        \multicolumn{1}{|l|}{Precondición}                           & \multicolumn{7}{p{10cm}|}{{\color[HTML]{808080} \textit{Un cliente (registrado o no) debe estar en el punto en el que tienen que pagar una reserva.}}}                                                                                                                                                                                                                                   \\ \hline
        \multicolumn{1}{|l|}{Poscondición}                           & \multicolumn{7}{p{10cm}|}{{\color[HTML]{808080} \textit{El dinero debe estar ingresado en la cuenta PayPal de la empresa.}}}                                                                                                                                                                                                                                                       \\ \hline
        \multicolumn{1}{|l|}{Autor}                                  & \multicolumn{1}{l|}{{\color[HTML]{808080} \textit{Jesús Rodríguez González}}}                           & \multicolumn{2}{l|}{Fecha}                                                   & \multicolumn{1}{l|}{}                          & \multicolumn{2}{l|}{Versión}                                                      & \multicolumn{1}{l|}{1.0}                                                     \\ \hline
        \end{tabular}
    \end{table}

    \begin{table}[H]
        \begin{tabular}{llllllll}
    \hline
    \multicolumn{8}{|l|}{\cellcolor[HTML]{F2F2F2}Propósito}                                                                                                                                                                                                                                                                                                                                                                                           \\ \hline
    \multicolumn{8}{|p{15cm}|}{{\color[HTML]{808080} \textit{Permitir a un cliente o cliente registrado efectuar el pago de una reserva mediante PayPal.}}}                                                                                                                                                                                                             \\ \hline
    \end{tabular}
    \end{table}

    \begin{table}[H]
        \begin{tabular}{llllllll}
    \hline
    \multicolumn{8}{|l|}{\cellcolor[HTML]{F2F2F2}Resumen}                                                                                                                                                                                                                                                                                                                                                                                             \\ \hline
    \multicolumn{8}{|p{15cm}|}{{\color[HTML]{808080} \textit{Tras haber comenzado la reserva de una actividad, un cliente debe poder elegir la opción de realizar el pago a través de PayPal, llevándose a cabo el pago por el cliente y recibiendo el dinero la empresa.}}} \\ \hline
    \end{tabular}
    \end{table}

    \begin{table}[h!]
        \centering
        \begin{tabular}{|p{0.6cm}|p{6.5cm}|p{0.6cm}|p{6.5cm}|}
        \rowcolor{lightgray}
        \hline
        \multicolumn{4}{|l|}{Curso Normal (Básico)} \\
        \hline
        1 & \color[HTML]{808080}\textit{Un cliente se dispone a pagar con PayPal.} & 2 & \color[HTML]{808080}\textit{El sistema informa al sistema PayPal de una petición de pago.} \\
        \hline
        3 & \color[HTML]{808080}\textit{El sistema PayPal se comunica con el sistema para poder comenzar el proceso de pago.} & 4 & \color[HTML]{808080}\textit{El sistema traslada al sistema PayPal los datos del cliente y permite proceder con el pago.} \\
        \hline
        5 & \color[HTML]{808080}\textit{El sistema PayPal envía al sistema una ventana emergente en la que el cliente pueda introducir sus datos de pago.} & 6 & \color[HTML]{808080}\textit{El sistema reenvía esta ventana emergente al cliente.} \\
        \hline
        7 & \color[HTML]{808080}\textit{El cliente introduce sus datos y procede al pago.} & 8 & \color[HTML]{808080}\textit{El sistema notifica que el pago ha sido realizado correctamente.} \\
        \hline
        \end{tabular}
        \end{table}
        
        \begin{table}[h!]
            \centering
            \begin{tabular}{|p{0.6cm}|p{13.4cm}|}
            \rowcolor{lightgray}
            \hline
            \multicolumn{2}{|l|}{Cursos Alternos} \\
            \hline
            3a & \color[HTML]{808080}\textit{Si el sistema no permite proceder con el pago, se notifica al cliente y termina el caso de uso.} \\
            \hline
            7a & \color[HTML]{808080}\textit{Si los datos introducidos son incorrectos, se notifica al cliente y termina el caso de uso.} \\
            \hline
            \end{tabular}
            \end{table}

            \begin{table}[H]
                \begin{tabular}{|l
                >{\columncolor[HTML]{FFFFFF}}l l
                >{\columncolor[HTML]{FFFFFF}}l 
                >{\columncolor[HTML]{FFFFFF}}l 
                >{\columncolor[HTML]{FFFFFF}}l 
                >{\columncolor[HTML]{FFFFFF}}l 
                >{\columncolor[HTML]{FFFFFF}}l |}
                \hline
                \multicolumn{8}{|l|}{\cellcolor[HTML]{F2F2F2}Otros datos}                                                                                                                                                                                                                                                                                                        \\ \hline
                \multicolumn{1}{|l|}{Frecuencia esperada} & \multicolumn{1}{p{5cm}|}{\cellcolor[HTML]{FFFFFF}{\color[HTML]{808080} }} & \multicolumn{1}{l|}{Rendimiento} & \multicolumn{5}{p{5cm}|}{\cellcolor[HTML]{FFFFFF}{\color[HTML]{808080} }} \\ \hline
                \multicolumn{1}{|l|}{Importancia}         & \multicolumn{1}{p{5cm}|}{\cellcolor[HTML]{FFFFFF}{\color[HTML]{808080} }}   & \multicolumn{1}{l|}{Urgencia}    & \multicolumn{5}{p{5cm}|}{\cellcolor[HTML]{FFFFFF}{\color[HTML]{808080} }}                          \\ \hline
                \multicolumn{1}{|l|}{Estado}              & \multicolumn{1}{p{5cm}|}{\cellcolor[HTML]{FFFFFF}{\color[HTML]{808080}}}                   & \multicolumn{1}{l|}{Estabilidad} & \multicolumn{5}{p{5cm}|}{\cellcolor[HTML]{FFFFFF}{\color[HTML]{808080} }}                                            \\ \hline
                \end{tabular}
                \end{table}

   
                \begin{table}[h!]
                    \centering
                    \begin{tabular}{|p{14.2cm}|}
                    \rowcolor{lightgray}
                    \hline
                    Comentarios \\
                    \hline
                     \\
                    \hline
                    \end{tabular}
                    \end{table}

    % ======Tabla Caso de uso Número: 6 =========================
    \begin{table}[H]
        \begin{tabular}{llllllll}
        \hline
        \multicolumn{1}{|l|}{Caso de Uso}                            & \multicolumn{5}{l|}{{\color[HTML]{808080} \textit{Pago con Bizum}}}                                                                                                                                                                & \multicolumn{2}{l|}{\cellcolor[HTML]{E6E6E6}{\color[HTML]{808080} \textit{CU-6}}}                                                   \\ \hline
        \multicolumn{1}{|l|}{Actores}                                & \multicolumn{7}{l|}{{\color[HTML]{808080} \textit{Cliente(I), Cliente Registrado, Banco}}}                                                                                                                                                                                                                                                                                               \\ \hline
        \multicolumn{1}{|l|}{Tipo}                                   & \multicolumn{7}{l|}{{\color[HTML]{808080} \textit{Primario Esencial}}}                                                                                                                                                                                                                                                                                                             \\ \hline
        \multicolumn{1}{|l|}{Referencias}                            & \multicolumn{2}{l|}{}                                                                                                   & \multicolumn{5}{l|}{\color[HTML]{808080}\textit{CU-4 Pagar reserva (padre).}}                                                                                                                                                                                                                                    \\ \hline
        \multicolumn{1}{|l|}{Precondición}                           & \multicolumn{7}{p{10cm}|}{{\color[HTML]{808080} \textit{Un cliente (registrado o no) debe estar en el punto en el que tienen que pagar una reserva.}}}                                                                                                                                                                                                                                   \\ \hline
        \multicolumn{1}{|l|}{Poscondición}                           & \multicolumn{7}{p{10cm}|}{{\color[HTML]{808080} \textit{El dinero debe estar ingresado en la cuenta bancaria de la empresa.}}}                                                                                                                                                                                                                                                       \\ \hline
        \multicolumn{1}{|l|}{Autor}                                  & \multicolumn{1}{l|}{{\color[HTML]{808080} \textit{Jesús Rodríguez González}}}                           & \multicolumn{2}{l|}{Fecha}                                                   & \multicolumn{1}{l|}{}                          & \multicolumn{2}{l|}{Versión}                                                      & \multicolumn{1}{l|}{1.0}                                                     \\ \hline
        \end{tabular}
    \end{table}

    \begin{table}[H]
        \begin{tabular}{llllllll}
    \hline
    \multicolumn{8}{|l|}{\cellcolor[HTML]{F2F2F2}Propósito}                                                                                                                                                                                                                                                                                                                                                                                           \\ \hline
    \multicolumn{8}{|p{15cm}|}{{\color[HTML]{808080} \textit{Permitir a un cliente o cliente registrado efectuar el pago de una reserva mediante Bizum.}}}                                                                                                                                                                                                             \\ \hline
    \end{tabular}
    \end{table}

    \begin{table}[H]
        \begin{tabular}{llllllll}
    \hline
    \multicolumn{8}{|l|}{\cellcolor[HTML]{F2F2F2}Resumen}                                                                                                                                                                                                                                                                                                                                                                                             \\ \hline
    \multicolumn{8}{|p{15cm}|}{{\color[HTML]{808080} \textit{Tras haber comenzado la reserva de una actividad, un cliente debe poder elegir la opción de realizar el pago a través de Bizum, llevándose a cabo el pago por el cliente y recibiendo el dinero la empresa.}}} \\ \hline
    \end{tabular}
    \end{table}
    \begin{table}[h!]
        \centering
        \begin{tabular}{|p{0.6cm}|p{6.5cm}|p{0.6cm}|p{6.5cm}|}
        \rowcolor{lightgray}
        \hline
        \multicolumn{4}{|l|}{Curso Normal (Básico)} \\
        \hline
        1 & \color[HTML]{808080}\textit{Un cliente se dispone a pagar con Bizum.} & 2 & \color[HTML]{808080}\textit{El sistema informa al banco de una petición de pago mediante Bizum.} \\
        \hline
        3 & \color[HTML]{808080}\textit{El banco se comunica con el sistema para poder comenzar el proceso de pago.} & 4 & \color[HTML]{808080}\textit{El sistema traslada al banco los datos del cliente y permite proceder con el pago.} \\
        \hline
        5 & \color[HTML]{808080}\textit{El banco envía al sistema una ventana emergente para que el cliente pueda introducir sus datos de pago con Bizum.} & 6 & \color[HTML]{808080}\textit{El sistema reenvía esta ventana emergente al cliente.} \\
        \hline
        7 & \color[HTML]{808080}\textit{El cliente introduce sus datos y procede al pago.} & 8 & \color[HTML]{808080}\textit{El sistema notifica que el pago ha sido realizado correctamente}. \\
        \hline
        \end{tabular}
        \end{table}
        

        \begin{table}[h!]
            \centering
            \begin{tabular}{|p{0.6cm}|p{13.4cm}|}
            \rowcolor{lightgray}
            \hline
            \multicolumn{2}{|l|}{Cursos Alternos} \\
            \hline
            3a & \color[HTML]{808080}\textit{Si el sistema no permite proceder con el pago, se notifica al cliente y termina el caso de uso.} \\
            \hline
            7a & \color[HTML]{808080}\textit{Si los datos introducidos son incorrectos, se notifica al cliente y termina el caso de uso.} \\
            \hline
            \end{tabular}
            \end{table}
            
            \begin{table}[H]
                \begin{tabular}{|l
                >{\columncolor[HTML]{FFFFFF}}l l
                >{\columncolor[HTML]{FFFFFF}}l 
                >{\columncolor[HTML]{FFFFFF}}l 
                >{\columncolor[HTML]{FFFFFF}}l 
                >{\columncolor[HTML]{FFFFFF}}l 
                >{\columncolor[HTML]{FFFFFF}}l |}
                \hline
                \multicolumn{8}{|l|}{\cellcolor[HTML]{F2F2F2}Otros datos}                                                                                                                                                                                                                                                                                                        \\ \hline
                \multicolumn{1}{|l|}{Frecuencia esperada} & \multicolumn{1}{p{5cm}|}{\cellcolor[HTML]{FFFFFF}{\color[HTML]{808080} }} & \multicolumn{1}{l|}{Rendimiento} & \multicolumn{5}{p{5cm}|}{\cellcolor[HTML]{FFFFFF}{\color[HTML]{808080} }} \\ \hline
                \multicolumn{1}{|l|}{Importancia}         & \multicolumn{1}{p{5cm}|}{\cellcolor[HTML]{FFFFFF}{\color[HTML]{808080} }}   & \multicolumn{1}{l|}{Urgencia}    & \multicolumn{5}{p{5cm}|}{\cellcolor[HTML]{FFFFFF}{\color[HTML]{808080} }}                          \\ \hline
                \multicolumn{1}{|l|}{Estado}              & \multicolumn{1}{p{5cm}|}{\cellcolor[HTML]{FFFFFF}{\color[HTML]{808080}}}                   & \multicolumn{1}{l|}{Estabilidad} & \multicolumn{5}{p{5cm}|}{\cellcolor[HTML]{FFFFFF}{\color[HTML]{808080} }}                                            \\ \hline
                \end{tabular}
                \end{table}


                \begin{table}[h!]
                    \centering
                    \begin{tabular}{|p{14.2cm}|}
                    \rowcolor{lightgray}
                    \hline
                    Comentarios \\
                    \hline
                     \\
                    \hline
                    \end{tabular}
                    \end{table}

   % ======Tabla Caso de uso Número: 7 =========================
   \begin{table}[H]
    \begin{tabular}{llllllll}
    \hline
    \multicolumn{1}{|l|}{Caso de Uso}                            & \multicolumn{5}{l|}{{\color[HTML]{808080} \textit{Pago con tarjeta}}}                                                                                                                                                                & \multicolumn{2}{l|}{\cellcolor[HTML]{E6E6E6}{\color[HTML]{808080} \textit{CU-7}}}                                                   \\ \hline
    \multicolumn{1}{|l|}{Actores}                                & \multicolumn{7}{l|}{{\color[HTML]{808080} \textit{Cliente(I), Cliente Registrado, Banco}}}                                                                                                                                                                                                                                                                                               \\ \hline
    \multicolumn{1}{|l|}{Tipo}                                   & \multicolumn{7}{l|}{{\color[HTML]{808080} \textit{Primario Esencial}}}                                                                                                                                                                                                                                                                                                             \\ \hline
    \multicolumn{1}{|l|}{Referencias}                            & \multicolumn{2}{l|}{}                                                                                                   & \multicolumn{5}{l|}{\color[HTML]{808080}\textit{CU-4 Pagar reserva (padre).}}                                                                                                                                                                                                                                    \\ \hline
    \multicolumn{1}{|l|}{Precondición}                           & \multicolumn{7}{p{10cm}|}{{\color[HTML]{808080} \textit{Un cliente (registrado o no) debe estar en el punto en el que tienen que pagar una reserva.}}}                                                                                                                                                                                                                                   \\ \hline
    \multicolumn{1}{|l|}{Poscondición}                           & \multicolumn{7}{p{10cm}|}{{\color[HTML]{808080} \textit{El dinero debe estar ingresado en la cuenta bancaria de la empresa.}}}                                                                                                                                                                                                                                                       \\ \hline
    \multicolumn{1}{|l|}{Autor}                                  & \multicolumn{1}{l|}{{\color[HTML]{808080} \textit{Jesús Rodríguez González}}}                           & \multicolumn{2}{l|}{Fecha}                                                   & \multicolumn{1}{l|}{}                          & \multicolumn{2}{l|}{Versión}                                                      & \multicolumn{1}{l|}{1.0}                                                     \\ \hline
    \end{tabular}
\end{table}

\begin{table}[H]
    \begin{tabular}{llllllll}
\hline
\multicolumn{8}{|l|}{\cellcolor[HTML]{F2F2F2}Propósito}                                                                                                                                                                                                                                                                                                                                                                                           \\ \hline
\multicolumn{8}{|p{15cm}|}{{\color[HTML]{808080} \textit{Permitir a un cliente o cliente registrado efectuar el pago de una reserva mediante tarjeta bancaria.}}}                                                                                                                                                                                                             \\ \hline
\end{tabular}
\end{table}

\begin{table}[H]
    \begin{tabular}{llllllll}
\hline
\multicolumn{8}{|l|}{\cellcolor[HTML]{F2F2F2}Resumen}                                                                                                                                                                                                                                                                                                                                                                                             \\ \hline
\multicolumn{8}{|p{15cm}|}{{\color[HTML]{808080} \textit{Tras haber comenzado la reserva de una actividad, un cliente debe poder elegir la opción de realizar el pago a través de tarjeta bancaria, llevándose a cabo el pago por el cliente y recibiendo el dinero la empresa.}}} \\ \hline
\end{tabular}
\end{table}

\begin{table}[h!]
    \centering
    \begin{tabular}{|p{0.6cm}|p{6.5cm}|p{0.6cm}|p{6.5cm}|}
    \rowcolor{lightgray}
    \hline
    \multicolumn{4}{|l|}{Curso Normal (Básico)} \\
    \hline
    1 & \color[HTML]{808080}\textit{Un cliente se dispone a pagar con tarjeta.} & 2 & \color[HTML]{808080}\textit{El sistema informa al banco de una petición de pago mediante tarjeta bancaria.} \\
    \hline
    3 & \color[HTML]{808080}\textit{El banco se comunica con el sistema para poder comenzar el proceso de pago.} & 4 & \color[HTML]{808080}\textit{El sistema traslada al banco los datos del cliente y permite proceder al pago.} \\
    \hline
    5 & \color[HTML]{808080}\textit{El banco envía al sistema una ventana emergente para que el cliente pueda introducir sus datos de pago con tarjeta.} & 6 & \color[HTML]{808080}\textit{El sistema reenvía esta ventana emergente al cliente.} \\
    \hline
    7 & \color[HTML]{808080}\textit{El cliente introduce sus datos y procede al pago.} & 8 & \color[HTML]{808080}\textit{El sistema notifica que ha sido realizado correctamente.} \\
    \hline
    \end{tabular}
    \end{table}

    \begin{table}[h!]
        \centering
        \begin{tabular}{|p{0.6cm}|p{13.4cm}|}
        \rowcolor{lightgray}
        \hline
        \multicolumn{2}{|l|}{Cursos Alternos} \\
        \hline
        3a & \color[HTML]{808080}\textit{Si el sistema no permite proceder con el pago, se notifica al cliente y termina el caso de uso.} \\
        \hline
        7a & \color[HTML]{808080}\textit{Si los datos introducidos son incorrectos, se notifica al cliente y termina el caso de uso.} \\
        \hline
        \end{tabular}
        \end{table}
        
        \begin{table}[H]
            \begin{tabular}{|l
            >{\columncolor[HTML]{FFFFFF}}l l
            >{\columncolor[HTML]{FFFFFF}}l 
            >{\columncolor[HTML]{FFFFFF}}l 
            >{\columncolor[HTML]{FFFFFF}}l 
            >{\columncolor[HTML]{FFFFFF}}l 
            >{\columncolor[HTML]{FFFFFF}}l |}
            \hline
            \multicolumn{8}{|l|}{\cellcolor[HTML]{F2F2F2}Otros datos}                                                                                                                                                                                                                                                                                                        \\ \hline
            \multicolumn{1}{|l|}{Frecuencia esperada} & \multicolumn{1}{p{5cm}|}{\cellcolor[HTML]{FFFFFF}{\color[HTML]{808080} }} & \multicolumn{1}{l|}{Rendimiento} & \multicolumn{5}{p{5cm}|}{\cellcolor[HTML]{FFFFFF}{\color[HTML]{808080} }} \\ \hline
            \multicolumn{1}{|l|}{Importancia}         & \multicolumn{1}{p{5cm}|}{\cellcolor[HTML]{FFFFFF}{\color[HTML]{808080} }}   & \multicolumn{1}{l|}{Urgencia}    & \multicolumn{5}{p{5cm}|}{\cellcolor[HTML]{FFFFFF}{\color[HTML]{808080} }}                          \\ \hline
            \multicolumn{1}{|l|}{Estado}              & \multicolumn{1}{p{5cm}|}{\cellcolor[HTML]{FFFFFF}{\color[HTML]{808080}}}                   & \multicolumn{1}{l|}{Estabilidad} & \multicolumn{5}{p{5cm}|}{\cellcolor[HTML]{FFFFFF}{\color[HTML]{808080} }}                                            \\ \hline
            \end{tabular}
            \end{table}


            \begin{table}[h!]
                \centering
                \begin{tabular}{|p{14.2cm}|}
                \rowcolor{lightgray}
                \hline
                Comentarios \\
                \hline
                 \\
                \hline
                \end{tabular}
                \end{table}

    % ======Tabla Caso de uso Número: 8 =========================
    \begin{table}[H]
        \begin{tabular}{llllllll}
        \hline
        \multicolumn{1}{|l|}{Caso de Uso}                            & \multicolumn{5}{l|}{{\color[HTML]{808080} \textit{Cancelar reserva}}}                                                                                                                                                                & \multicolumn{2}{l|}{\cellcolor[HTML]{E6E6E6}{\color[HTML]{808080} \textit{CU-8}}}                                                   \\ \hline
        \multicolumn{1}{|l|}{Actores}                                & \multicolumn{7}{l|}{{\color[HTML]{808080} \textit{Cliente Regsitrado (I)}}}                                                                                                                                                                                                                                                                                               \\ \hline
        \multicolumn{1}{|l|}{Tipo}                                   & \multicolumn{7}{l|}{{\color[HTML]{808080} \textit{Primario Esencial}}}                                                                                                                                                                                                                                                                                                             \\ \hline
        \multicolumn{1}{|l|}{Referencias}                            & \multicolumn{2}{l|}{}                                                                                                   & \multicolumn{5}{l|}{}                                                                                                                                                                                                                                    \\ \hline
        \multicolumn{1}{|l|}{Precondición}                           & \multicolumn{7}{p{10cm}|}{{\color[HTML]{808080} \textit{El cliente debe estar registrado y tener una reserva previa, además de estar dentro de los plazos preestablecidos para poder cancelarla.}}}                                                                                                                                                                                                                                   \\ \hline
        \multicolumn{1}{|l|}{Poscondición}                           & \multicolumn{7}{p{10cm}|}{{\color[HTML]{808080} \textit{Al cliente se le ha devuelto el dinero pagado y ya no tiene la reserva.}}}                                                                                                                                                                                                                                                       \\ \hline
        \multicolumn{1}{|l|}{Autor}                                  & \multicolumn{1}{l|}{{\color[HTML]{808080} \textit{Jesús Rodríguez González}}}                           & \multicolumn{2}{l|}{Fecha}                                                   & \multicolumn{1}{l|}{}                          & \multicolumn{2}{l|}{Versión}                                                      & \multicolumn{1}{l|}{1.0}                                                     \\ \hline
        \end{tabular}
    \end{table}

    \begin{table}[H]
        \begin{tabular}{llllllll}
        \hline
        \multicolumn{8}{|l|}{\cellcolor[HTML]{F2F2F2}Propósito}                                                                                                                                                                                                                                                                                                                                                                                           \\ \hline
        \multicolumn{8}{|p{15cm}|}{{\color[HTML]{808080} \textit{Permitir a un cliente registrado llevar a cabo la cancelación de una reserva.}}}                                                                                                                                                                                                             \\ \hline
        \end{tabular}
    \end{table}

    \begin{table}[H]
        \begin{tabular}{llllllll}
        \hline
        \multicolumn{8}{|l|}{\cellcolor[HTML]{F2F2F2}Resumen}                                                                                                                                                                                                                                                                                                                                                                                             \\ \hline
        \multicolumn{8}{|p{15cm}|}{{\color[HTML]{808080} \textit{Si un cliente ha hecho una reserva estando registrado y ses encuentra dentro de unos plazos de tiempo previamente establecidos, podrá llevar a cabo una cancelación de su reserva siendo devuelto el dinero abonado para la misma.}}} \\ \hline
        \end{tabular}
    \end{table}
    
    \begin{table}[h!]
        \centering
        \begin{tabular}{|p{0.6cm}|p{6.5cm}|p{0.6cm}|p{6.5cm}|}
        \rowcolor{lightgray}
        \hline
        \multicolumn{4}{|l|}{Curso Normal (Básico)} \\
        \hline
        1 & \color[HTML]{808080}\textit{Un cliente registrado se dispone a cancelar una reserva.} & 2 & \color[HTML]{808080}\textit{El sistema verifica los datos del cliente registrado.} \\
        \hline
        & & 3 & \color[HTML]{808080}\textit{El sistema verifica si la reserva que desea cancelar está almacenada en su base de datos.} \\
        \hline
        & & 4 & \color[HTML]{808080}\textit{El sistema verifica si se cumplen los plazos establecidos para poder cancelar una reserva.} \\
        \hline
        & & 5 & \color[HTML]{808080}\textit{Se realiza el reembolso de la reserva, incluir CU-10 (Reembolso reserva)}. \\
        \hline
        \end{tabular}
        \end{table}
        
    \begin{table}[h!]
\centering
\begin{tabular}{|p{0.6cm}|p{13.4cm}|}
\rowcolor{lightgray}
\hline
\multicolumn{2}{|l|}{Cursos Alternos} \\
\hline
2a & \color[HTML]{808080}\textit{Si datos del cliente no son válidos, se notifica al cliente y termina el caso de uso.} \\
\hline
3a & \color[HTML]{808080}\textit{Si no se encuentra la reserva a cancelar, se notifica al cliente y termina el caso de uso.} \\
\hline
4a & \color[HTML]{808080}\textit{Si los plazos no se cumplen, se notifica al cliente y termina el caso de uso.} \\
\hline
\end{tabular}
\end{table}


\begin{table}[H]
    \begin{tabular}{|l
    >{\columncolor[HTML]{FFFFFF}}l l
    >{\columncolor[HTML]{FFFFFF}}l 
    >{\columncolor[HTML]{FFFFFF}}l 
    >{\columncolor[HTML]{FFFFFF}}l 
    >{\columncolor[HTML]{FFFFFF}}l 
    >{\columncolor[HTML]{FFFFFF}}l |}
    \hline
    \multicolumn{8}{|l|}{\cellcolor[HTML]{F2F2F2}Otros datos}                                                                                                                                                                                                                                                                                                        \\ \hline
    \multicolumn{1}{|l|}{Frecuencia esperada} & \multicolumn{1}{p{5cm}|}{\cellcolor[HTML]{FFFFFF}{\color[HTML]{808080} }} & \multicolumn{1}{l|}{Rendimiento} & \multicolumn{5}{p{5cm}|}{\cellcolor[HTML]{FFFFFF}{\color[HTML]{808080} }} \\ \hline
    \multicolumn{1}{|l|}{Importancia}         & \multicolumn{1}{p{5cm}|}{\cellcolor[HTML]{FFFFFF}{\color[HTML]{808080} }}   & \multicolumn{1}{l|}{Urgencia}    & \multicolumn{5}{p{5cm}|}{\cellcolor[HTML]{FFFFFF}{\color[HTML]{808080} }}                          \\ \hline
    \multicolumn{1}{|l|}{Estado}              & \multicolumn{1}{p{5cm}|}{\cellcolor[HTML]{FFFFFF}{\color[HTML]{808080}}}                   & \multicolumn{1}{l|}{Estabilidad} & \multicolumn{5}{p{5cm}|}{\cellcolor[HTML]{FFFFFF}{\color[HTML]{808080} }}                                            \\ \hline
    \end{tabular}
    \end{table}


    \begin{table}[h!]
        \centering
        \begin{tabular}{|p{14.2cm}|}
        \rowcolor{lightgray}
        \hline
        Comentarios \\
        \hline
         \\
        \hline
        \end{tabular}
        \end{table}

    % ======Tabla Caso de uso Número: 9 =========================
    \begin{table}[H]
        \begin{tabular}{llllllll}
        \hline
        \multicolumn{1}{|l|}{Caso de Uso}                            & \multicolumn{5}{l|}{{\color[HTML]{808080} \textit{Ver mis reservas}}}                                                                                                                                                                & \multicolumn{2}{l|}{\cellcolor[HTML]{E6E6E6}{\color[HTML]{808080} \textit{CU-9}}}                                                   \\ \hline
        \multicolumn{1}{|l|}{Actores}                                & \multicolumn{7}{l|}{{\color[HTML]{808080} \textit{Cliente Regsitrado(I)}}}                                                                                                                                                                                                                                                                                               \\ \hline
        \multicolumn{1}{|l|}{Tipo}                                   & \multicolumn{7}{l|}{{\color[HTML]{808080} \textit{Primario Esencial}}}                                                                                                                                                                                                                                                                                                             \\ \hline
        \multicolumn{1}{|l|}{Referencias}                            & \multicolumn{2}{l|}{}                                                                                                   & \multicolumn{5}{l|}{\color[HTML]{808080}\textit{CU-10 Reembolso reserva (incluye)}}                                                                                                                                                                                                                                    \\ \hline
        \multicolumn{1}{|l|}{Precondición}                           & \multicolumn{7}{p{10cm}|}{{\color[HTML]{808080} \textit{}}}                                                                                                                                                                                                                                   \\ \hline
        \multicolumn{1}{|l|}{Poscondición}                           & \multicolumn{7}{p{10cm}|}{{\color[HTML]{808080} \textit{}}}                                                                                                                                                                                                                                                       \\ \hline
        \multicolumn{1}{|l|}{Autor}                                  & \multicolumn{1}{l|}{{\color[HTML]{808080} \textit{Jesús Rodríguez González}}}                           & \multicolumn{2}{l|}{Fecha}                                                   & \multicolumn{1}{l|}{}                          & \multicolumn{2}{l|}{Versión}                                                      & \multicolumn{1}{l|}{1.0}                                                     \\ \hline
        \end{tabular}
    \end{table}

    \begin{table}[H]
        \begin{tabular}{llllllll}
        \hline
        \multicolumn{8}{|l|}{\cellcolor[HTML]{F2F2F2}Propósito}                                                                                                                                                                                                                                                                                                                                                                                           \\ \hline
        \multicolumn{8}{|p{15cm}|}{{\color[HTML]{808080} \textit{Permite que un cliente registrado pueda ver información sobre las reservas que ha realizado con la aplicación.}}}                                                                                                                                                                                                             \\ \hline
        \end{tabular}
    \end{table}

    \begin{table}[H]
        \begin{tabular}{llllllll}
        \hline
        \multicolumn{8}{|l|}{\cellcolor[HTML]{F2F2F2}Resumen}                                                                                                                                                                                                                                                                                                                                                                                             \\ \hline
        \multicolumn{8}{|p{15cm}|}{{\color[HTML]{808080} \textit{Un cliente que haya hecho una serie de reservas estando registrado con la misma cuenta puede ver el historial de todas las resevas que ha hecho a lo largo del tiempo a través de la aplicación.}}} \\ \hline
        \end{tabular}
    \end{table}

    \begin{table}[h!]
        \centering
        \begin{tabular}{|p{0.6cm}|p{6.5cm}|p{0.6cm}|p{6.5cm}|}
        \rowcolor{lightgray}
        \hline
        \multicolumn{4}{|l|}{Curso Normal (Básico)} \\
        \hline
        1 & \color[HTML]{808080}\textit{Un cliente registrado solicita ver sus reservas.} & 2 & \color[HTML]{808080}\textit{El sistema comprueba que el usuario está registrado.} \\
        \hline
        & & 3 & \color[HTML]{808080}\textit{El sistema busca las reservas del cliente almacenadas en su base de datos.} \\
        \hline
        & & 4 & \color[HTML]{808080}\textit{El sistema muestra al cliente la lista de sus reservas almacenada.} \\
        \hline
        \end{tabular}
        \end{table}
        

        \begin{table}[h!]
            \centering
            \begin{tabular}{|p{0.6cm}|p{13.4cm}|}
            \rowcolor{lightgray}
            \hline
            \multicolumn{2}{|l|}{Cursos Alternos} \\
            \hline
            2a & \color[HTML]{808080}\textit{Si el sistema comprueba que el cliente no está registrado, se notifica al cliente y termina el caso de uso.} \\
            \hline
            \end{tabular}
            \end{table}
            

            \begin{table}[H]
                \begin{tabular}{|l
                >{\columncolor[HTML]{FFFFFF}}l l
                >{\columncolor[HTML]{FFFFFF}}l 
                >{\columncolor[HTML]{FFFFFF}}l 
                >{\columncolor[HTML]{FFFFFF}}l 
                >{\columncolor[HTML]{FFFFFF}}l 
                >{\columncolor[HTML]{FFFFFF}}l |}
                \hline
                \multicolumn{8}{|l|}{\cellcolor[HTML]{F2F2F2}Otros datos}                                                                                                                                                                                                                                                                                                        \\ \hline
                \multicolumn{1}{|l|}{Frecuencia esperada} & \multicolumn{1}{p{5cm}|}{\cellcolor[HTML]{FFFFFF}{\color[HTML]{808080} }} & \multicolumn{1}{l|}{Rendimiento} & \multicolumn{5}{p{5cm}|}{\cellcolor[HTML]{FFFFFF}{\color[HTML]{808080} }} \\ \hline
                \multicolumn{1}{|l|}{Importancia}         & \multicolumn{1}{p{5cm}|}{\cellcolor[HTML]{FFFFFF}{\color[HTML]{808080} }}   & \multicolumn{1}{l|}{Urgencia}    & \multicolumn{5}{p{5cm}|}{\cellcolor[HTML]{FFFFFF}{\color[HTML]{808080} }}                          \\ \hline
                \multicolumn{1}{|l|}{Estado}              & \multicolumn{1}{p{5cm}|}{\cellcolor[HTML]{FFFFFF}{\color[HTML]{808080}}}                   & \multicolumn{1}{l|}{Estabilidad} & \multicolumn{5}{p{5cm}|}{\cellcolor[HTML]{FFFFFF}{\color[HTML]{808080} }}                                            \\ \hline
                \end{tabular}
                \end{table}
    
                \begin{table}[h!]
                    \centering
                    \begin{tabular}{|p{14.2cm}|}
                    \rowcolor{lightgray}
                    \hline
                    Comentarios \\
                    \hline
                     \\
                    \hline
                    \end{tabular}
                    \end{table}
   